\section{合卷：战争怎样重做国家}

十九世纪中叶的战争没有简单地毁掉国家，也没有直接制造一个整齐的“地方化”国家。清廷在省级军队、厘金、捐输、外交机构、企业和赈济网络中获得继续作战和重建的资源；同一过程也让税源、武装和信息越来越依赖各省的中介者。

太平、捻军、西北、西南和对外战争的时间并不相同。把它们放在同一财政与道路网络中，才能看见一地的河工、港口、难民或军械为何会改变另一地的选择。战后企业与新机构不是脱离战争的“自强”插曲，而是对军费、航运、外交和人才短缺的回应。

\subsection{制度得失}

\textbf{所得：} 省级筹饷和新式机构增加了应对危机的资源，外交与技术事务有了较稳定的处理渠道。

\textbf{代价：} 地方军事、财政和精英网络成为不可绕开的政治条件；战争损失、族群暴力和难民重建无法由胜利报告结案。

\textbf{留下的压力：} 国家若要改革，必须面对既有省份网络，而不能只靠宫廷发布新的制度名称。
